% \clearpage
% \pagestyle{empty}
%
% \addcontentsline{toc}{chapter}{PEDOMAN PENGGUNAAN TESIS}
%
% \begin{center}
%     \fontsize{14}{17}\selectfont \bfseries \MakeUppercase{Pedoman Penggunaan Tesis}
% \end{center}
%
% Tesis Magister yang tidak dipublikasikan terdaftar dan tersedia di Perpustakaan
% Institut Teknologi Bandung, dan terbuka untuk umum dengan ketentuan bahwa hak
% cipta ada pada penulis dengan mengikuti aturan HaKI yang berlaku di Institut
% Teknologi Bandung. Referensi kepustakaan diperkenankan dicatat, tetapi
% pengutipan atau peringkasan hanya dapat dilakukan seizin penulis dan harus
% disertai dengan kaidah ilmiah untuk menyebutkan sumbernya.
%
% Sitasi hasil penelitian Tesis ini dapat di tulis dalam bahasa Indonesia sebagai berikut:
% \vspace{-1.5\baselineskip}
% \begin{list}{}{\setlength{\leftmargin}{1.27cm} \setlength{\itemindent}{-\leftmargin}}
% \item Nama Belakang, Inisial Nama Depan. (Tahun): \textit{Judul Tesis}, Tesis Program Magister, Institut Teknologi Bandung.
% \end{list}
%
% dan dalam bahasa Inggris sebagai berikut:
%
% \begin{list}{}{\setlength{\leftmargin}{1.27cm} \setlength{\itemindent}{-\leftmargin}}
% \item Nama Belakang, Inisial Nama Depan. (Tahun): \textit{Judul tesis yang telah diterjemahkan dalam bahasa Inggris}, Master's Thesis, Institut Teknologi Bandung.
% \end{list}
%
% Memperbanyak atau menerbitkan sebagian atau seluruh tesis haruslah seizin Dekan
% Sekolah Pascasarjana, Institut Teknologi Bandung.
%
% \clearpage